\section{Persistencia durable en Cloud Storage}

La durabilidad del mini-banco se apoya en un \emph{log} de transferencias que
vive en un bucket de Google Cloud Storage. Cada transferencia confirmada por el
lider termina escrita en esa bitacora, de modo que si caen los tres nodos a la
vez el lider pueda renacer y reconstruir todos los saldos releyendo el log. Tres
clases cooperan: \codigo{durable/Journal.java} orquesta la escritura y la
recuperacion, \codigo{durable/GcsStore.java} habla con la JSON API de Cloud
Storage y \codigo{durable/GcsAuth.java} obtiene los tokens de acceso. La unidad
que se persiste es siempre un \codigo{core/Transfer.java}, el mismo registro que
viaja por el cluster.

\subsection{La bitacora de transacciones}

El \codigo{durable/Journal.java} es un orquestador de durabilidad que solo existe
en el lider (las replicas nunca poseen un \codigo{Journal}). Implementa la
interfaz \codigo{core/CommitListener.java}: una vez registrado sobre el banco del
lider, su metodo \codigo{onCommit(Transfer)} recibe cada transferencia recien
confirmada. Lo importante es que \emph{no} sube nada a GCS en ese momento: solo
encola la transferencia y regresa de inmediato. Asi, la respuesta HTTP de una
transferencia nunca se bloquea esperando un viaje de ida y vuelta a Cloud
Storage.

\begin{lstlisting}[language=Java,caption={onCommit solo encola; la subida ocurre fuera del hot path}]
@Override
public void onCommit(Transfer transfer) {
    queue.offer(transfer);
}
\end{lstlisting}

La cola es un \codigo{BlockingQueue} que un unico hilo de fondo
(\codigo{journal-writer}, marcado como \emph{daemon}) consume en
\codigo{drainLoop()}. Ese hilo espera con \codigo{poll(POLL\_MS, ...)} a que
aparezca una transferencia, y cuando la toma vacia de golpe todo lo que haya
disponible en la cola hasta un maximo de \codigo{MAX\_BATCH} (1000) con
\codigo{drainTo}, componiendo un solo lote. El lote completo se sube como un
unico objeto mediante \codigo{store.putBatch(batch)} y, solo si la subida tuvo
exito, se avanza el contador \codigo{stored} en el tamano del lote. El agrupado
por lotes es lo que permite que el log durable siga el ritmo de una tasa alta de
commits sin un objeto por transferencia.

\begin{lstlisting}[language=Java,caption={drainLoop compone un lote y lo sube como un solo objeto}]
batch.clear();
batch.add(first);
queue.drainTo(batch, MAX_BATCH - 1);
try {
    store.putBatch(batch);
    stored.addAndGet(batch.size());
} catch (IOException e) {
    System.err.println("journal: could not persist batch of " + batch.size()
            + " ending at seq " + batch.get(batch.size() - 1).seq()
            + ": " + e.getMessage());
}
\end{lstlisting}

Cada entrada del log es exactamente un \codigo{core/Transfer.java}: un registro
inmutable con un numero de secuencia monotono asignado por el ledger al confirmar
(\codigo{seq}), la cuenta origen (\codigo{from}), la cuenta destino
(\codigo{to}) y el monto movido en centavos (\codigo{cents}). Es la misma forma
que serializan el codec del cluster y el journal de GCS.

\begin{lstlisting}[language=Java,caption={La unidad persistida: el registro Transfer (core/Transfer.java)}]
public record Transfer(long seq, int from, int to, long cents) {
}
\end{lstlisting}

El contador \codigo{stored()} reporta cuantas transferencias estan
\emph{realmente} subidas a GCS; queda por debajo del numero de commits aplicados
en exactamente la profundidad de la cola, que es la cifra durable que debe
mostrar el dashboard. Un fallo de subida se registra en \codigo{stderr} pero no
detiene el hilo: el lazo corre mientras \codigo{running} sea verdadero o quede
algo en la cola. Un lote cuya subida falla se descarta y deja un hueco en el log
durable; para que ese hueco no pase inadvertido, \codigo{recover()} valida la
contiguidad al arrancar (el log es contiguo desde 1) y avisa por \codigo{stderr} si
una secuencia no es la sucesora esperada, en vez de recuperar en silencio un log
incompleto.

Esta estrategia asincrona tiene una ventana de durabilidad explicita. En un
apagado ordenado, \codigo{stop()} pone \codigo{running} en falso y hace
\codigo{join} al escritor para drenar la cola, de modo que toda transferencia ya
confirmada quede subida antes de que el proceso termine; \codigo{app/Node.java}
registra ese \codigo{stop()} como \emph{shutdown hook}. En cambio, una muerte
abrupta del proceso deja la cola en vuelo como la ventana de durabilidad inherente
a un journaling asincrono: como el commit responde al cliente antes de subir a
GCS, una transferencia confirmada puede perderse en un \emph{kill} duro. Ese es
el precio consciente de mantener Cloud Storage fuera del camino critico.

\subsection{Acceso a Cloud Storage}

El \codigo{durable/GcsStore.java} es un cliente delgado sobre la JSON API REST de
Cloud Storage, construido directamente con \codigo{java.net.http.HttpClient} y
sin el SDK de Google. Todos los objetos del journal cuelgan del prefijo
\codigo{journal/}. El metodo \codigo{put(Transfer)} existe y escribiria una
transferencia suelta como \codigo{journal/tx-seq.json}, pero el camino vivo no lo
usa: \codigo{Journal} siempre llama a \codigo{putBatch}, que compone muchas
transferencias en un solo objeto \codigo{journal/batch-min-max.json} (un arreglo
JSON), nombrado por la secuencia minima y maxima del lote. Cloud Storage no tiene
\emph{append}; componer objetos es la alternativa sancionada por el contrato.

Las subidas usan el endpoint de \emph{simple media upload}
(\codigo{UPLOAD\_BASE}) con \codigo{uploadType=media}, mientras que listar y leer
usan la JSON API (\codigo{API\_BASE}). Toda peticion lleva la cabecera
\codigo{Authorization: Bearer} con un token de \codigo{GcsAuth} y se serializa con
gson. Ademas, \codigo{GcsStore} ofrece \codigo{putCheckpoint}/\codigo{getCheckpoint}
para que una replica guarde su snapshot bajo el prefijo \codigo{checkpoint/};
estos quedan deliberadamente fuera de \codigo{journal/} para que el
\codigo{readAll()} del lider nunca los liste ni intente parsearlos como
\codigo{Transfer}.

\begin{lstlisting}[language=Java,caption={putBatch compone el lote en un solo objeto journal/batch-min-max.json}]
long min = batch.stream().mapToLong(Transfer::seq).min().getAsLong();
long max = batch.stream().mapToLong(Transfer::seq).max().getAsLong();
String name = PREFIX + "batch-" + min + "-" + max + ".json";
String url = UPLOAD_BASE + bucket + "/o?uploadType=media&name=" + enc(name);
HttpRequest request = HttpRequest.newBuilder(URI.create(url))
        .timeout(HTTP_TIMEOUT)
        .header("Authorization", "Bearer " + auth.accessToken())
        .header("Content-Type", "application/json")
        .POST(HttpRequest.BodyPublishers.ofString(gson.toJson(batch)))
        .build();
\end{lstlisting}

Toda peticion pasa por \codigo{sendWithRetry}, que reintenta hasta
\codigo{MAX\_ATTEMPTS} (3) veces ante un fallo de red transitorio
(\codigo{IOException}: una conexion caida, un timeout) y, si todos los intentos
fallan, relanza la excepcion. Esto hace que la recuperacion en frio aborte de
forma ruidosa en lugar de servir silenciosamente un estado solo-base. Un estado
HTTP no-2xx no se reintenta: se devuelve al llamador, que decide como tratarlo
(por ejemplo, \codigo{getCheckpoint} interpreta un 404 como ``aun no existe'' y
devuelve \codigo{null}).

\subsubsection{Autenticacion con la cuenta de servicio}

El \codigo{durable/GcsAuth.java} acuna tokens de acceso OAuth2 para la JSON API
usando el flujo de \emph{key-file} de cuenta de servicio, tambien sin SDK de
Google (el metadata server de la VM no se usa a proposito). En el constructor lee
el JSON de la cuenta de servicio apuntado por \codigo{TES\_GCS\_KEYFILE} y extrae
\codigo{client\_email}, \codigo{token\_uri} (con un endpoint por defecto si falta)
y \codigo{private\_key}, del que reconstruye la llave RSA PKCS8.

Para obtener un token, \codigo{buildAssertion()} arma un JWT con cabecera RS256 y
un claim set cargado con \codigo{iss}, \codigo{scope}
(\codigo{devstorage.read\_write}), \codigo{aud}, \codigo{iat} y \codigo{exp}, lo
firma con \codigo{SHA256withRSA} sobre la llave privada y lo codifica en Base64
URL sin relleno. No se envia \codigo{sub} (eso es solo para delegacion de
dominio y aqui seria rechazado como \codigo{unauthorized\_client}).

\begin{lstlisting}[language=Java,caption={Claims del JWT que GcsAuth autofirma (RS256)}]
JsonObject claims = new JsonObject();
claims.addProperty("iss", clientEmail);
claims.addProperty("scope", SCOPE);
claims.addProperty("aud", tokenUri);
claims.addProperty("iat", now);
claims.addProperty("exp", now + TOKEN_TTL_SECONDS);
\end{lstlisting}

El JWT firmado se intercambia en \codigo{refresh()} por un POST
\codigo{x-www-form-urlencoded} al \codigo{token\_uri} con el grant
\codigo{urn:ietf:params:oauth:grant-type:jwt-bearer}, y del cuerpo se lee
\codigo{access\_token}. \codigo{accessToken()} cachea el token y lo reusa hasta
poco antes de expirar (margen \codigo{REFRESH\_SKEW\_SECONDS} de 60 s), acunando
uno nuevo solo cuando hace falta; el metodo es \codigo{synchronized} porque el
hilo escritor lo consulta en cada subida.

El alta de la cuenta de servicio se hace una sola vez, segun describe
\codigo{SETUP-GCS.md}: crear la cuenta, otorgarle \codigo{roles/storage.objectAdmin}
sobre el bucket (no basta \codigo{objectCreator}, porque el log no solo sube sino
que \emph{relee} en el arranque en frio) y descargar la llave
\codigo{credentials.json}. Esa llave nunca se sube al repositorio; en la VM se
copia aparte y se apunta \codigo{TES\_GCS\_KEYFILE} a su ruta, mientras
\codigo{TES\_BUCKET} lleva el nombre del bucket (solo en el nodo lider).

\subsection{Recuperacion ante desastre total}

Cuando el lider arranca, \codigo{app/Node.java} corre la recuperacion durable
antes de aceptar trafico, siempre que Cloud Storage este configurado (si faltan
\codigo{TES\_BUCKET} o \codigo{TES\_GCS\_KEYFILE}, el lider sigue sirviendo todos
los endpoints en memoria y el journal simplemente no se activa). El arranque crea
el \codigo{GcsAuth} y el \codigo{GcsStore}, instancia el \codigo{Journal} y llama
a \codigo{recover}, pasandole \codigo{bank::applyReplicated} como la accion que
aplica cada transferencia recuperada al ledger.

\begin{lstlisting}[language=Java,caption={Arranque del lider: recuperar, sembrar el contador y registrar el journal (app/Node.java)}]
GcsAuth auth = new GcsAuth(Path.of(config.gcsKeyfile()));
GcsStore store = new GcsStore(auth, config.bucket());
journal = new Journal(store);
int recovered = journal.recover(bank::applyReplicated);
journal.start(recovered);
bank.addCommitListener(journal);
\end{lstlisting}

Dentro de \codigo{recover(CommitListener apply)}, el journal pide a
\codigo{store.readAll()} todas las transferencias persistidas y reaplica cada una
en orden al ledger. El metodo \codigo{readAll()} de \codigo{GcsStore} lista los
objetos bajo el prefijo \codigo{journal/} (siguiendo la paginacion por
\codigo{nextPageToken} de la JSON API), baja cada uno y los aplana: si el objeto
es un arreglo (un lote) expande cada elemento, y si es un objeto suelto lo
convierte directo. Como cada \codigo{Transfer} carga su \codigo{seq}, el ``donde
quedo'' no depende del nombre de los objetos ni de un cursor persistido del
lider: \codigo{readAll()} reordena todo por secuencia al final, asi que el ledger
se reconstruye exactamente hasta la ultima secuencia confirmada sin huecos ni
duplicados de orden.

\begin{lstlisting}[language=Java,caption={readAll aplana lotes y transferencias sueltas, y reordena por secuencia (GcsStore)}]
JsonElement parsed = JsonParser.parseString(response.body());
if (parsed.isJsonArray()) {
    for (JsonElement element : parsed.getAsJsonArray()) {
        transfers.add(gson.fromJson(element, Transfer.class));
    }
} else {
    transfers.add(gson.fromJson(parsed, Transfer.class));
}
// ...
transfers.sort(Comparator.comparingLong(Transfer::seq));
\end{lstlisting}

Hecha la reaplicacion, \codigo{recover} devuelve cuantas transferencias se
recuperaron y \codigo{start(recovered)} siembra con ese valor el contador
\codigo{stored} (para que el dashboard arranque mostrando lo ya durable) y lanza
el hilo escritor que sube los commits vivos a partir de ese punto. El lider queda
entonces sirviendo con los saldos completos reconstruidos desde el log, que el
dashboard muestra como saldos recuperados. La garantia es fuerte para todo lo que
alcanzo a subirse a GCS y se vio reforzada por el \emph{shutdown hook} que drena
la cola en un apagado ordenado; lo unico no recuperable es la ventana asincrona
en vuelo descrita antes, ante una caida abrupta.
